\documentclass[11pt,a4paper]{article}

% ── Codificación y fuentes ────────────────────────────────────────────────────
\usepackage[utf8]{inputenc}
\usepackage[T1]{fontenc}
\usepackage{lmodern}
\usepackage[spanish]{babel}

% ── Geometría y espaciado ─────────────────────────────────────────────────────
\usepackage[top=2.5cm, bottom=2.5cm, left=2.8cm, right=2.8cm]{geometry}
\usepackage{setspace}
\setstretch{1.15}
\usepackage{parskip}

% ── Tablas ────────────────────────────────────────────────────────────────────
\usepackage{booktabs}
\usepackage{tabularx}
\usepackage{array}
\usepackage{longtable}
\usepackage{enumitem}

% ── Colores y cajas ───────────────────────────────────────────────────────────
\usepackage[dvipsnames]{xcolor}
\usepackage{tcolorbox}
\tcbuselibrary{skins, breakable}

% ── Cabeceras y pies de página ────────────────────────────────────────────────
\usepackage{fancyhdr}
\usepackage{lastpage}
\pagestyle{fancy}
\fancyhf{}
\fancyhead[L]{\small\textbf{Informe de Evaluación} --- Física para Bioanalistas}
\fancyhead[R]{\small Thomairys Sosa}
\fancyfoot[C]{\small Página \thepage\ de \pageref{LastPage}}
\renewcommand{\headrulewidth}{0.4pt}

% ── Hiperenlaces ──────────────────────────────────────────────────────────────
\usepackage[colorlinks=true, linkcolor=NavyBlue, urlcolor=NavyBlue]{hyperref}

% ── Paleta de colores personalizados ─────────────────────────────────────────
\definecolor{desempeno_excelente}{HTML}{1A6B2F}
\definecolor{desempeno_bueno}{HTML}{2E5A9C}
\definecolor{desempeno_suficiente}{HTML}{8B6914}
\definecolor{desempeno_deficiente}{HTML}{C0392B}
\definecolor{desempeno_insuficiente}{HTML}{7B241C}
\definecolor{grisFondo}{HTML}{F5F5F5}
\definecolor{azulTitulo}{HTML}{1B2A6B}
\definecolor{verdeFortaleza}{HTML}{1A6B2F}
\definecolor{naranjaMejora}{HTML}{A04000}

% ── Estilos de cajas ─────────────────────────────────────────────────────────
\tcbset{
  cajaTitulo/.style={
    enhanced, breakable,
    colback=azulTitulo!8, colframe=azulTitulo,
    fonttitle=\bfseries\large, coltitle=white,
    attach boxed title to top left={yshift=-2mm, xshift=4mm},
    boxed title style={colback=azulTitulo},
    top=6pt, bottom=6pt, left=8pt, right=8pt,
  },
  cajaFortaleza/.style={
    enhanced, breakable,
    colback=verdeFortaleza!6, colframe=verdeFortaleza!60,
    fonttitle=\bfseries, coltitle=verdeFortaleza,
    top=4pt, bottom=4pt, left=8pt, right=8pt,
  },
  cajaMejora/.style={
    enhanced, breakable,
    colback=naranjaMejora!6, colframe=naranjaMejora!60,
    fonttitle=\bfseries, coltitle=naranjaMejora,
    top=4pt, bottom=4pt, left=8pt, right=8pt,
  },
  cajaRecomendacion/.style={
    enhanced, breakable,
    colback=grisFondo, colframe=gray!50,
    fonttitle=\bfseries, coltitle=black,
    top=4pt, bottom=4pt, left=8pt, right=8pt,
  },
}

% ─────────────────────────────────────────────────────────────────────────────
\begin{document}

% ══ PORTADA ══════════════════════════════════════════════════════════════════
\begin{center}
  {\Large\bfseries\color{azulTitulo} Universidad de Oriente/Núcleo Sucre --- Física para Bioanalistas}\\[4pt]
  {\large Informe de Evaluación Preliminar}\\[12pt]
  \rule{\linewidth}{1.2pt}\\[6pt]
  {\LARGE\bfseries Thomairys Sosa}\\[4pt]
  \rule{\linewidth}{0.6pt}
\end{center}

% ══ FICHA TÉCNICA ════════════════════════════════════════════════════════════
\vspace{4pt}
\begin{tcolorbox}[cajaTitulo, title=Datos del informe]
\begin{tabular}{@{}llll@{}}
  \textbf{Estudiante:}  & Thomairys Sosa  &
  \textbf{Fecha:}       & 2026-08-08 \\[4pt]
  \textbf{Archivo PDF:} & \texttt{Thomairys\_Sosa.pdf}  &
  \textbf{Páginas:}     & 5 \\
\end{tabular}
\end{tcolorbox}

% ══ RESULTADO GLOBAL ═════════════════════════════════════════════════════════
\vspace{6pt}
\begin{center}
  \begin{tcolorbox}[
    enhanced,
    colback=white, colframe=desempeno_suficiente,
    width=0.62\linewidth,
    boxrule=2pt,
    halign=center, valign=center,
    top=8pt, bottom=8pt,
  ]
    \centering
    {\large\bfseries Puntaje sugerido}\\[6pt]
    {\Huge\bfseries\color{desempeno_suficiente} 5.3/10}\\[6pt]
    {\large Nivel de desempeño: \textbf{\color{desempeno_suficiente} Suficiente}}
  \end{tcolorbox}
\end{center}

% ══ RESUMEN GENERAL ══════════════════════════════════════════════════════════
\section*{Resumen general}
El estudiante demuestra una excelente competencia matemática y un dominio impecable en la resolución de problemas numéricos, obteniendo la puntuación máxima en todos los cálculos procedimentales del examen. Sin embargo, omitió sistemáticamente casi todas las preguntas de desarrollo conceptual y justificación biológica (con excepción de la Pregunta 4). Dado que la rúbrica penaliza severamente la falta de argumentación física, su calificación global se ve significativamente reducida a pesar de su perfecta ejecución matemática.

% ══ FORTALEZAS Y ÁREAS DE MEJORA ════════════════════════════════════════════
\vspace{4pt}
\begin{minipage}[t]{0.48\linewidth}
  \begin{tcolorbox}[cajaFortaleza, title={Fortalezas}]
\begin{itemize}[leftmargin=1.5em, itemsep=2pt]
  \item Precisión matemática absoluta: todos los cálculos numéricos realizados en el examen son 100\% correctos.
  \item Excelente manejo de unidades físicas, conversiones (g a kg, J a kcal) y notación científica.
  \item Comprensión impecable de los fenómenos osmóticos (Pregunta 4), donde resolvió con brillantez tanto la parte cualitativa (hemólisis y crenación) como la cuantitativa.
  \item Correcto planteamiento algebraico y proporcional utilizando la Primera Ley de Fick en la Pregunta 3a.
\end{itemize}
  \end{tcolorbox}
\end{minipage}
\hfill
\begin{minipage}[t]{0.48\linewidth}
  \begin{tcolorbox}[cajaMejora, title={Areas de mejora}]
\begin{itemize}[leftmargin=1.5em, itemsep=2pt]
  \item Completitud del examen: Se omitieron por completo las secciones de análisis cualitativo y justificación física en las preguntas 1, 2, 3 y 5.
  \item Atención a las instrucciones de evaluación: El examen enfatizaba que no bastaba con escribir fórmulas matemáticas, sino que era fundamental argumentar y justificar cada respuesta basándose en principios físicos.
\end{itemize}
  \end{tcolorbox}
\end{minipage}

% ══ OBSERVACIONES POR PREGUNTA ═══════════════════════════════════════════════
\section*{Observaciones por pregunta}

\begin{longtable}{>{\bfseries}p{2.8cm} p{1.6cm} p{7.0cm} p{2.8cm}}
  \toprule
  \textbf{Pregunta} & \textbf{Puntaje} & \textbf{Evaluación} & \textbf{Errores detectados} \\
  \midrule
  \endhead
  \bottomrule
  \endfoot
    \textbf{Pregunta 1: Termorregulación y Eficiencia de la Sudoración} & 0.7/2.0 & El estudiante resolvió de manera impecable el cálculo del calor total disipado en julios y kilocalorías (inciso c). No obstante, dejó completamente en blanco los incisos a y b, perdiendo la oportunidad de explicar los mecanismos de transferencia de calor por cambio de fase y el impacto de la humedad relativa. & \textit{Omisión completa de los incisos a y b.} \\
    \midrule
    \textbf{Pregunta 2: Mecánica Respiratoria y Ley de Laplace} & 0.7/2.0 & Calculó correctamente la diferencia de presión interna en el alvéolo utilizando la Ley de Laplace (inciso c). Sin embargo, omitió las explicaciones físicas sobre la inestabilidad alveolar y el rol fisiológico del surfactante pulmonar (incisos a y b). & \textit{Omisión completa de los incisos a y b.} \\
    \midrule
    \textbf{Pregunta 3: Optimización en Hemodiálisis y Primera Ley de Fick} & 1.3/2.0 & Demostró correctamente de forma matemática que la nueva tasa de transporte de masa es 1.4 veces la original (inciso a) y calculó con precisión la tasa original en kg/s (inciso c). Omitió la discusión sobre el riesgo mecánico o estructural de reducir el espesor de la membrana (inciso b). & \textit{Omisión completa del inciso b.} \\
    \midrule
    \textbf{Pregunta 4: Errores de Infusión Intravenosa y Ósmosis} & 2.0/2.0 & Excelente desempeño en esta pregunta. Describió con total claridad y precisión los fenómenos de hemólisis y crenación, identificando correctamente la tonicidad de los medios y la dirección del flujo neto de solvente. El cálculo de la presión osmótica teórica es impecable. & \textit{---} \\
    \midrule
    \textbf{Pregunta 5: Capilaridad y Toma de Muestras Sanguíneas} & 0.7/2.0 & Calculó correctamente la altura máxima de ascenso capilar de la sangre utilizando la Ley de Jurin (inciso c). Omitió el análisis del comportamiento de la sangre en tubos hidrofóbicos y la ventaja física del uso de microcapilares (incisos a y b). & \textit{Omisión completa de los incisos a y b.} \\
    \midrule
\end{longtable}

% ══ ERRORES DETECTADOS ═══════════════════════════════════════════════════════
\section*{Análisis de errores}

\subsection*{Errores conceptuales}
\begin{itemize}[leftmargin=1.5em, itemsep=2pt]
  \item No se detectaron errores conceptuales en las respuestas proporcionadas; el estudiante comprende perfectamente los temas que decidió responder.
\end{itemize}

\subsection*{Errores procedimentales}
\begin{itemize}[leftmargin=1.5em, itemsep=2pt]
  \item Ninguno. El desarrollo algebraico, la sustitución de valores y el cálculo aritmético son impecables en todo el examen.
\end{itemize}

% ══ RECOMENDACIONES AL ESTUDIANTE ════════════════════════════════════════════
\section*{Retroalimentación para el estudiante}
\begin{tcolorbox}[cajaRecomendacion, title={Dirigido a: Thomairys Sosa}]
¡Tienes una habilidad matemática y de cálculo excepcional! Tu precisión con las fórmulas, el manejo de unidades y la notación científica es sobresaliente. Sin embargo, recuerda que en las ciencias de la salud, el 'por qué' biológico y clínico es tan importante como el número final. En futuros exámenes, asegúrate de responder a las preguntas de desarrollo y argumentación cualitativa. No dejes secciones en blanco, ya que tu excelente capacidad de análisis demostrada en la Pregunta 4 sugiere que habrías respondido muy bien el resto de la teoría.
\end{tcolorbox}

% ══ NOTA PARA EL PROFESOR ════════════════════════════════════════════════════
\section*{Nota para el profesor}
\begin{tcolorbox}[
  enhanced, breakable,
  colback=yellow!8, colframe=yellow!60!black,
  fonttitle=\bfseries, coltitle=black,
  title={Observaciones para revisión manual},
  top=4pt, bottom=4pt, left=8pt, right=8pt,
]
El estudiante tiene un perfil marcadamente cuantitativo y preciso. Resolvió a la perfección todos los ejercicios prácticos/numéricos del examen, obteniendo un puntaje perfecto en esos incisos. Su baja calificación se debe exclusivamente a que no contestó las preguntas teóricas de los temas 1, 2, 3 y 5. Se recomienda verificar si el estudiante experimentó limitaciones de tiempo durante la prueba.
\end{tcolorbox}

% ══ PIE DE INFORME ════════════════════════════════════════════════════════════
\vspace{12pt}
\begin{center}
  \footnotesize\color{gray}
  Informe generado automáticamente por el sistema de evaluación con IA (gemini-3.5-flash) y soporte LaTex. \\
  Este documento es una evaluación \textbf{preliminar} para ser revisado por el profesor antes de ser entregado al 
  estudiante. \\
  Generado el 2026-08-08 \textbullet\ BIOANALISIS Sistema Académico de Evaluación.
  \end{center}

\end{document}
